% Généré par tools/build.py — ne pas éditer à la main
\section{Droites du plan}

\fichehead{lun.~2~nov.~2026}{Vecteur directeur d'une droite}
\begin{fiche}
\qitem{Échauffement · Flash}{GEO}{Donner un vecteur colinéaire à $\vec u\vc{2}{-4}$ à coordonnées entières plus simples.}{$\vc{1}{-2}$ (on divise par $2$).}{}
\qitem{Cœur · Méthode}{GEO}{Donner un vecteur directeur de $(AB)$ pour $A(1;2)$ et $B(4;8)$.}{$\vect{AB}\vc{3}{6}$, ou $\vc{1}{2}$ qui lui est colinéaire.}{}
\qitem{Cœur + · QCM}{GEO}{Un vecteur directeur de la droite passant par $A(0;1)$ et $B(2;-3)$ est~:\choix{$\vc{1}{-2}$}{$\vc{1}{-4}$}{$\vc{-4}{-2}$}{$\vc{2}{4}$}}{A.\ $\vc{1}{-2}$}{Ici $\vect{AB}\vc{2}{-4}$\,; on teste la colinéarité par le déterminant.}
\qitem{Retour · Flash}{STA}{Notes~: $8,\,12,\,12,\,15,\,18$. Calculer la moyenne, la médiane et l'étendue.}{Moyenne $13$\,; médiane $12$\,; étendue $10$.}{}
\qitem{Fil rouge · Propriété}{DEM}{$ABC$ est rectangle en $A$ et $I$ est le milieu de $[BC]$. Que peut-on dire de $IA$, $IB$ et $IC$\,? Pourquoi\,?}{$IA=IB=IC$~: le milieu de l'hypoténuse est le centre du cercle circonscrit.}{}
\end{fiche}

\fichehead{mar.~3~nov.~2026}{Équation cartésienne d'une droite}
\begin{fiche}
\qitem{Échauffement · Flash}{SUB}{Calculer $3x-2y+1$ pour $x=1$ et $y=-2$.}{$3+4+1=8$.}{}
\qitem{Cœur · Méthode}{GEO}{La droite $d$ a pour équation $2x+3y-6=0$. Donner un vecteur directeur et un point de $d$.}{Pour $ax+by+c=0$, $\vc{-b}{a}$ dirige la droite~: $\vc{-3}{2}$. Le point $(3;0)$ convient~: $6+0-6=0$.}{}
\qitem{Cœur + · Vrai/faux}{GEO}{Vrai ou faux\,? Justifier. Les droites $2x+3y-6=0$ et $4x+6y+1=0$ sont parallèles.}{Vrai. Vecteurs directeurs $\vc{-3}{2}$ et $\vc{-6}{4}$ colinéaires\,; elles sont distinctes~: $(3;0)$ n'appartient pas à la seconde droite ($12+1\neq0$).}{}
\qitem{Retour · Tableau}{LEC}{Mode de transport de $30$ élèves.\par\smallskip\renewcommand{\arraystretch}{1.15}\begin{tabular}{|c|c|c|c|c|}\hline
Transport & Bus & Vélo & Marche & Voiture \\ \hline
Effectif & 12 & 5 & 8 & 5 \\ \hline
\end{tabular}\par Donner la part des élèves venant en bus et l'angle du secteur «~vélo~» d'un diagramme circulaire.}{$\dfrac{12}{30}=40\,\%$\,; $\dfrac{5}{30}\times360\degre=60\degre$.}{}
\qitem{Fil rouge · Flash}{REP}{Sur une droite graduée, $A$ a pour abscisse $-2$ et $B$ pour abscisse $3$. Donner $AB$ et l'abscisse du milieu de $[AB]$.}{$AB=3-(-2)=5$\,; milieu~: $\dfrac{-2+3}{2}=0{,}5$.}{}
\end{fiche}

\fichehead{jeu.~5~nov.~2026}{Équation réduite et pente d'une droite non verticale}
\begin{fiche}
\qitem{Échauffement · Flash}{EQ}{Isoler $y$ dans $2x+y-5=0$.}{$y=-2x+5$.}{}
\qitem{Cœur · Méthode}{GEO}{Donner la pente et l'ordonnée à l'origine de la droite $3x-2y+4=0$.}{$-2y=-3x-4$, $y=\dfrac32x+2$~: pente $\dfrac32$, ordonnée à l'origine $2$.}{Piège~: diviser par $-2$ chaque terme (signes).}
\qitem{Cœur + · QCM}{GEO}{La pente de la droite $(AB)$ avec $A(1;2)$ et $B(5;0)$ est~:\choix{$2$}{$\dfrac12$}{$-2$}{$-\dfrac12$}}{D.\ $-\dfrac12$}{Pente $=\dfrac{y_B-y_A}{x_B-x_A}=\dfrac{-2}{4}$.}
\qitem{Retour · Flash}{CAL}{Calculer $\dfrac34+\dfrac56$.}{$\dfrac{9}{12}+\dfrac{10}{12}=\dfrac{19}{12}$.}{}
\qitem{Fil rouge · Problème}{DIV}{Deux engrenages ont $12$ et $18$ dents. Après combien de tours de chacun reviennent-ils ensemble à leur position de départ\,?}{$\ppcm(12\,;\,18)=36$ dents~: $3$ tours du premier, $2$ tours du second.}{}
\end{fiche}

\fichehead{lun.~9~nov.~2026}{Tracer une droite donnée par une équation}
\begin{fiche}
\qitem{Échauffement · Flash}{IMG}{Pour $y=-2x+3$, calculer $y$ lorsque $x=0$ puis lorsque $x=2$.}{$y=3$ puis $y=-1$.}{}
\qitem{Cœur · Méthode}{REG}{Donner deux points permettant de tracer la droite $2x+y-4=0$.}{Pour $x=0$~: $y=4$\,; pour $x=2$~: $y=0$. Points $(0;4)$ et $(2;0)$.}{}
\qitem{Cœur + · Flash}{GEO}{Quelle est la droite d'équation $x=3$\,? Donner un vecteur directeur. A-t-elle une pente\,?}{La verticale passant par $(3;0)$\,; $\vc{0}{1}$ la dirige\,; elle n'a pas de pente\,; équation cartésienne $x-3=0$.}{Ne s'écrit pas $y=mx+p$.}
\qitem{Retour · Flash}{CMP}{Comparer $\dfrac37$ et $\dfrac49$.}{Produits en croix~: $3\times9=27<28=4\times7$, donc $\dfrac37<\dfrac49$.}{}
\qitem{Fil rouge · Flash}{EQ}{Résoudre dans $\R$~: $5x-7=2x+8$.}{$x=5$}{Piège~: regrouper les $x$ d'un côté en changeant correctement de signe.}
\end{fiche}

\fichehead{mar.~10~nov.~2026}{Déterminer une équation de droite~: deux points, point et vecteur directeur, point et pente}
\begin{fiche}
\qitem{Échauffement · Flash}{GEO}{Pente de la droite passant par $(1;3)$ et $(3;7)$\,?}{$\dfrac{7-3}{3-1}=2$.}{}
\qitem{Cœur · Méthode}{GEO}{Équation réduite de la droite de pente $2$ passant par $A(1;3)$.}{$y-3=2(x-1)$, soit $y=2x+1$.}{Contrôle~: $2\times1+1=3$.}
\qitem{Cœur + · Méthode}{GEO}{Droite passant par $A(1;2)$ de vecteur directeur $\vec u\vc{2}{-1}$~: équation cartésienne par le déterminant de $\vect{AM}$ et $\vec u$.}{$M(x;y)\in d\Leftrightarrow\deter{\vect{AM}}{\vec u}=0$~: $(x-1)(-1)-(y-2)\times2=0$, soit $x+2y-5=0$.}{Contrôle~: $1+4-5=0$\,; $\vc{-2}{1}$ est bien colinéaire à $\vc{2}{-1}$.}
\qitem{Retour · Flash}{PUI}{Simplifier $2^5\times2^{-3}$.}{$2^{5-3}=2^2=4$.}{}
\qitem{Fil rouge · Tableau → formule}{REG}{Une fonction affine $f$ vérifie~: $f(0)=-1$, $f(1)=1$, $f(2)=3$. Donner $f(x)$.}{Pente $2$, ordonnée à l'origine $-1$~: $f(x)=2x-1$.}{}
\end{fiche}

\fichehead{jeu.~12~nov.~2026}{Droites et fonctions affines~: lien entre les registres}
\begin{fiche}
\qitem{Échauffement · Flash}{IMG}{$f(x)=-3x+4$. Calculer $f(2)$ et l'antécédent de $10$.}{$f(2)=-2$\,; $-3x+4=10\Leftrightarrow x=-2$.}{}
\qitem{Cœur · Méthode}{REG}{Déterminer la fonction affine $f$ telle que $f(1)=1$ et $f(3)=5$.}{Pente $\dfrac{5-1}{3-1}=2$\,; $f(x)=2x+p$ avec $2+p=1$~: $f(x)=2x-1$.}{}
\qitem{Cœur + · Vrai/faux}{REG}{Vrai ou faux\,? Justifier. Toute droite du plan est la courbe d'une fonction affine.}{Faux. Contre-exemple~: la droite verticale $x=3$ n'est la courbe d'aucune fonction (le nombre $3$ aurait une infinité d'images).}{}
\qitem{Retour · Flash}{PRB}{On lance un dé équilibré à $6$ faces. Probabilité d'obtenir un multiple de $3$\,?}{Issues $3$ et $6$~: $\dfrac26=\dfrac13$.}{}
\qitem{Fil rouge · Flash}{GEO}{$ABC$ est rectangle en $A$ avec $AB=6$ cm et $AC=8$ cm. Calculer $BC$.}{$BC^2=36+64=100$, donc $BC=10$ cm.}{}
\end{fiche}

\fichehead{lun.~16~nov.~2026}{Alignement, parallélisme, intersection de droites}
\begin{fiche}
\qitem{Échauffement · Flash}{SUB}{Le point $(3;5)$ appartient-il à la droite $y=2x-1$\,?}{$2\times3-1=5$~: oui.}{}
\qitem{Cœur · Méthode}{GEO}{Positions relatives de $d:\ y=3x+1$ avec $d':\ y=3x-2$, puis avec $d'':\ y=-x+1$.}{$d$ et $d'$~: même pente, ordonnées à l'origine distinctes~: parallèles strictement\,; $d$ et $d''$~: pentes différentes~: sécantes.}{}
\qitem{Cœur + · Méthode}{GEO}{Montrer que $A(1;4)$, $B(2;7)$ et $C(4;13)$ sont alignés de deux façons.}{(i) Tous vérifient $y=3x+1$. (ii) $\vect{AB}\vc{1}{3}$, $\vect{AC}\vc{3}{9}$~: $\deter{\vect{AB}}{\vect{AC}}=9-9=0$.}{}
\qitem{Retour · Flash}{PRO}{Sur une carte à l'échelle $\dfrac{1}{50\,000}$, deux villes sont distantes de $4$ cm. Quelle est la distance réelle\,?}{$4\times50\,000=200\,000$ cm $=2$ km.}{}
\qitem{Fil rouge · Méthode}{ALG}{Programme~: choisir un nombre $x$\,; ajouter $3$\,; multiplier par $2$\,; soustraire $6$. Calculer le résultat pour $x=5$, puis écrire une expression en $x$ et la réduire. Que remarque-t-on\,?}{Pour $5$~: $10$. Expression~: $2(x+3)-6=2x$~: on obtient toujours le double du nombre choisi.}{Lien A09~: l'écriture littérale prouve ce que les essais numériques suggèrent.}
\end{fiche}

\fichehead{mar.~17~nov.~2026}{Intersection de deux droites\,; fonction Python d'alignement}
\begin{fiche}
\qitem{Échauffement · Flash}{EQ}{Résoudre $2x+1=-x+7$.}{$3x=6$, donc $x=2$.}{}
\qitem{Cœur · Méthode}{GEO}{Coordonnées de l'intersection de $y=2x+1$ et $y=-x+7$.}{$x=2$ puis $y=2\times2+1=5$~: $(2;5)$.}{Contrôle dans la seconde équation~: $-2+7=5$.}
\qitem{Cœur + · Lecture de code}{ALG}{Que renvoient \code{alignes(0,0,1,2,2,4)} et \code{alignes(0,0,1,2,2,5)}\,?\par\smallskip{\ttfamily\small \hspace*{0.0em}def alignes(xa, ya, xb, yb, xc, yc):\\ \hspace*{2.4em}return (xb-xa)*(yc-ya) - (yb-ya)*(xc-xa) == 0\par}}{\code{True} puis \code{False}~: déterminants $1\times4-2\times2=0$ et $1\times5-2\times2=1\neq0$.}{}
\qitem{Retour · Développer}{DEV}{Développer et réduire $3(x-2)$.}{$3x-6$.}{}
\qitem{Fil rouge · Prouver}{LIT}{Montrer que la somme de trois entiers consécutifs est un multiple de $3$.}{Avec $n-1,\ n,\ n+1$~: somme $=3n$, multiple de $3$.}{Choisir le terme du milieu simplifie le calcul.}
\end{fiche}
